\section{Data and MC Samples}
\label{sec:monoZ_samples}

This search utilizes the full ATLAS Run 2 dataset ($\sqrt{s}=13~\text{TeV}$), corresponding to an integrated luminosity of $139~\text{fb}^{-1}$. 

MC simulations are used to model the detector response and signal acceptance for both the DM signal processes of interest alongside the SM backgrounds. The passage of particles through the ATLAS detector was simulated for all generated events relying on the \textsc{Geant4} toolkit~\cite{GEANT4:2002zbu} or a fast simulation where indicated, and reconstructed using the same software as the data. Pile-up effects (additional inelastic $pp$ collisions occurring within the identical or adjacent bunch crossings) were modelled by overlaying simulated minimum-bias events generated with \textsc{Pythia} 8 \cite{Sjostrand:2007gs}.

\subsection{Signal Samples}
\label{subsec:signal_samples}

Three distinct classes of signal models are simulated to interpret the results of this search: the invisible decay of the Higgs boson, simplified DM models with spin-1 mediators, and extended Two-Higgs-Doublet Models (2HDM+$a$).

\paragraph{\textbf{Higgs boson decay to invisible particles:}}
The SM $ZH$ production is generated utilizing the \textsc{Powheg Box} v2 framework \cite{Nason:2009xu} achieving Next-to-Leading-Order (NLO) accuracy in QCD. Both quark-induced ($q\bar{q} \to ZH$) and gluon-induced ($gg \to ZH$) processes are generated. The events are subsequently processed with \textsc{Pythia} 8 \cite{Sjostrand:2007gs} to model the parton showering, hadronisation, and underlying event kinematics, using the AZNLO tune and the NNPDF3.0 PDF set. The samples are normalised to cross-sections calculated at Next-to-Next-to-Leading Order (NNLO) in QCD with NLO Electroweak (EW) corrections \cite{deFlorian:2016spz}. The Higgs boson mass is set to $m_H = 125~\text{GeV}$, and it is required to decay to four neutrinos to mimic the invisible DM signature.

\paragraph{\textbf{Simplified Dark Matter models:}}
Signals for vector and axial-vector mediators decaying into Dirac Dark Matter particles ($\chi$) are generated using \textsc{MadGraph5\_aMC@NLO} v2.6.2 \cite{Alwall:2014hca} at NLO accuracy. The generation uses the \texttt{DMsimp\_s\_spin1} model implementation. Events are interfaced with \textsc{Pythia} 8 using the A14 set of tuned parameters (tune) and the NNPDF3.0 PDF set. The coupling of the mediator to the DM particles ($g_{\chi}$) is set to 1.0, and the universal coupling to quarks ($g_{q}$) is set to 0.25. A grid of samples is produced spanning the mediator mass ($m_{\text{med}}$) and DM mass ($m_{\chi}$) plane.

\paragraph{\textbf{2HDM+$a$ models:}}
The 2HDM+$a$ signal samples are generated using \textsc{MadGraph5\_aMC@NLO} \cite{Alwall:2014hca} at Leading Order (LO) accuracy. Both gluon-gluon fusion and $b$-quark associated production mechanisms are simulated. The subsequent phases of parton showering and hadronisation are executed via \textsc{Pythia} 8 with the A14 tune. The kinematics of the signal are scanned across four parameters: the pseudoscalar masses $m_a$ and $m_A$, the ratio of the vacuum expectation values corresponding to the two Higgs doublets, $\tan\beta$, as well as the mixing angle $\sin\theta$.

\subsection{Background Samples}
\label{subsec:bkg_samples}

The dominant background contributions arise from diboson ($ZZ, WZ, WW$) production and $Z+\text{jets}$ processes. Other contributions include top-quark pair production ($t\bar{t}$), single top, and triboson processes.

\paragraph{\textbf{Diboson ($ZZ, WZ, WW$):}}
The $q\bar{q}$-initiated diboson processes ($ZZ \to \ell\ell\nu\nu$, $ZZ \to 4\ell$, $WZ \to \ell\nu\ell'\ell'$, $WW \to \ell\nu \ell'\nu'$) are modeled by employing the \textsc{Sherpa} 2.2.2 generator \cite{Bothmann:2019yzt}. The matrix element calculations attain NLO precision for final states containing up to one extra parton, whereas LO precision is applied for up to three additional partons. The loop-induced process $gg \to ZZ$ is simulated at LO accuracy using \textsc{Sherpa} 2.2.2 with up to one additional parton. All diboson samples use the NNPDF3.0 NNLO PDF set.

\paragraph{\textbf{$Z+\text{jets}$:}}
Events comprising a $Z$ boson produced in conjunction with jets are simulated via \textsc{Sherpa} 2.2.1 \cite{Bothmann:2019yzt}. Matrix element evaluations achieve NLO accuracy for configurations with up to two partons and LO accuracy for those with up to four partons, using the Comix and OpenLoops libraries \cite{Cascioli:2011va}. The samples are normalised to a Next-to-Next-to-Leading Order (NNLO) prediction. This background is particularly relevant for events with fake $E_{\text{T}}^{\text{miss}}$ arising from jet mis-measurement.

\paragraph{\textbf{Top quark processes:}}
Top-quark pair ($t\bar{t}$) and single-top production ($Wt$, $s$-channel, and $t$-channel) are simulated using \textsc{Powheg Box} v2 \cite{Nason:2009xu, Hamilton:2012rf} interfaced with \textsc{Pythia} 8 \cite{Sjostrand:2007gs}. The A14 tune and NNPDF3.0 PDF set are used. The inclusive cross-sections for the $t\bar{t}$ samples are scaled to NNLO theoretical predictions including soft-gluon resummation to next-to-next-to-leading-log (NNLL) accuracy.

\paragraph{\textbf{Other Backgrounds:}}
Rare backgrounds such as triboson ($VVV$) and the associated production of top-quark pairs with a vector boson ($t\bar{t}V$), are simulated utilizing \textsc{Sherpa} 2.2.2 \cite{Bothmann:2019yzt} and \textsc{MadGraph5\_aMC@NLO} \cite{Alwall:2014hca}, respectively. While their contribution to the signal region is minor ($<1\%$), they are included for completeness in the background estimation.